\chapter{FUTURE WORK}
\label{ch:future-work}
Future studies could examine repeated exposure, task context, and the separate contributions of perceived resemblance and acoustic naturalness. Manipulating resemblance while controlling other acoustic properties would address a narrower causal question than the present comparison with a participant-selected neutral voice.

Other extensions could examine a visible agent body, a known-target search task, or a broader participant population. Each would change how the voice is encountered and would require a revised rationale and method.

Memory, provider autonomy, and gaze-guidance policy could also be varied in separate systems studies. Those comparisons concern the supporting mechanism and cannot be inferred from the present voice contrast. Likewise, eye-event or pupil analyses would be warranted only where validated measurements can address a defined question about task demands or voice experience.

\section{Gaze-Informed Adaptation of Spoken Assistance}
Nerella et al.'s gaze-based prediction study also motivates testing whether momentary task demands can inform when an agent speaks~\cite{nerella2026gaze}. For Follow My Voice, such an extension would require task-specific labels, evaluation on unseen participants, and a separate test of adaptive hint timing. The cited work leaves evaluation of real-time deployment open, and its classification result does not show that adapting spoken guidance improves search.

\section{Complementing Behavioral Measures with Brain Signals}
A future extension could combine electroencephalography (EEG) with gaze, task performance, and questionnaires to investigate the attentional demands of receiving self-similar and neutral spoken guidance. Yang et al.'s \emph{Synchronizing with Attentional Sampling} provides a relevant precedent: the authors use EEG in immersive VR to investigate attentional sampling during gaze and hand--eye interaction as part of a study of visual flicker guidance~\cite{yang2026attentional}. This work motivates examining neural activity alongside observable search behavior, although its visual-guidance findings do not establish how voice resemblance affects attention or how frequently spoken hints should occur.

For the present voice comparison, synchronized recordings could relate neural activity to actual instruction onset, orientation toward the target, and controller selection. A follow-up study could then examine whether the two voices are associated with different responses during these stages, even when completion times are similar. Such measurements would complement participants' reports of workload and experience rather than serve as direct measures of enjoyment, trust, or perceived voice similarity.

This extension would first require a feasibility study of EEG and headset compatibility, event synchronization, and the effects of eye movements, head turns, and muscle activity on the recordings. Neural measures and analysis windows would need to be specified in advance, with comparable task demands and speech exposure across voice conditions.
